\documentclass[11pt,a4paper]{article}
\usepackage[utf8]{inputenc}
\usepackage[spanish]{babel}
\usepackage{amsmath,amssymb,amsfonts}
\usepackage{geometry}
\usepackage{listings}
\usepackage{xcolor}

\geometry{margin=2.5cm}

\lstset{
    language=Maple,
    basicstyle=\ttfamily\small,
    breaklines=true,
    keywordstyle=\color{blue}\bfseries,
    commentstyle=\color{gray}\itshape,
    stringstyle=\color{red},
    frame=single,
    showstringspaces=false,
    backgroundcolor=\color{black!5}
}

\title{\textbf{Reporte Técnico: Resolución de Cálculo Vectorial con MAPLE}}
\author{Ing. Melvin E. Padilla}
\date{\today}

\begin{document}

\maketitle

\section{Introducción}
El presente documento detalla el desarrollo analítico y la implementación computacional en \textbf{MAPLE} para la resolución de problemas de cálculo vectorial aplicados a la ingeniería. Se abordan el modelado geométrico y cálculo de masa de un recipiente industrial, la verificación del Teorema de Stokes y la evaluación de integrales de línea mediante campos vectoriales conservativos.

\section{Problema 3: Geometría y Análisis de Sólidos}
\subsection{Parámetros Geométricos del Recipiente}
Se consideran las siguientes dimensiones para el recipiente de sección circular con fondo cónico:
\begin{itemize}
    \item Altura de la parte cilíndrica: $H = 10\text{ pie}$
    \item Altura de la parte cónica: $h = 10\text{ pie}$
    \item Diámetro de la carcasa: $D = \frac{2}{3}(H + h) = \frac{40}{3}\text{ pie}$
    \item Radio de la sección circular: $R = \frac{20}{3}\text{ pie}$
\end{itemize}

\subsection{Inciso a: Cálculo de la Masa}
La densidad de masa volumétrica en cada punto es proporcional al cuadrado de la distancia al punto más bajo del sólido:
$$\rho(r, z) = k(r^2 + z^2)$$
Empleando coordenadas cilíndricas, la masa total se obtiene integrando la densidad sobre el dominio correspondiente al fondo cónico y a la carcasa cilíndrica.

\subsection{Inciso b: Verificación del Teorema de Stokes}
Dado el campo vectorial $\mathbf{F}(x,y,z) = (x,y,z)$, su rotacional es idénticamente nulo ($\nabla \times \mathbf{F} = \mathbf{0}$), lo que permite verificar de forma directa la igualdad establecida por el Teorema de Stokes sobre la superficie de la cabeza del recipiente.

\section{Problema 4: Integral de Línea y Función Potencial}
Se evalúa la integral de línea del campo vectorial:
$$\mathbf{F}(x,y,z) = \left(4xe^z, \cos(y), 2x^2e^z\right)$$
A partir de la condición de campo conservativo, se determina la función potencial $f(x,y,z)$ y se aplica el Teorema Fundamental del Cálculo para trayectorias poligonales tridimensionales.

\section{Script de MAPLE Implementado}
A continuación se muestra el código utilizado en MAPLE para automatizar el cálculo analítico y las representaciones gráficas espaciales:

\begin{lstlisting}
restart:
with(plots):
with(VectorCalculus):
SetCoordinates(cartesian[x, y, z]):

# 1. VISUALIZACIÓN 3D DEL RECIPIENTE (PROBLEMA 3a)
H := 10:
h_cono := 10:
D := (2/3)*(H + h_cono):
R := D / 2:

p_cono := plot3d([(R/h_cono)*z*cos(theta), (R/h_cono)*z*sin(theta), z], 
                 z = 0 .. h_cono, theta = 0 .. 2*Pi, 
                 grid = [30, 30], color = Navy, shading = Z, 
                 transparency = 0.1, axes = boxed, 
                 title = "Geometria Completa del Recipiente Industrial"):

p_cilindro := plot3d([R*cos(theta), R*sin(theta), z], 
                     z = h_cono .. h_cono + H, theta = 0 .. 2*Pi, 
                     grid = [30, 30], color = Cyan, shading = Z, 
                     transparency = 0.2):

p_cabeza := plot3d([R*sin(phi)*cos(theta), R*sin(phi)*sin(theta), (h_cono + H) + R*cos(phi)], 
                   phi = 0 .. Pi/2, theta = 0 .. 2*Pi, 
                   grid = [20, 30], color = SteelBlue, 
                   transparency = 0.1):

display([p_cono, p_cilindro, p_cabeza], 
        scaling = constrained, 
        labels = ["X (pie)", "Y (pie)", "Z (pie)"],
        view = [-10..10, -10..10, -2..25]);

# 2. CÁLCULO DE MASA (PROBLEMA 3a)
rho := k * (r^2 + z^2):
dV := r:
M_cono := int(int(int(rho * dV, r = 0 .. (R/h_cono)*z), z = 0 .. h_cono), theta = 0 .. 2*Pi):
M_cilindro := int(int(int(rho * dV, r = 0 .. R), z = h_cono .. h_cono + H), theta = 0 .. 2*Pi):
M_total := simplify(M_cono + M_cilindro);

# 3. TEOREMA DE STOKES (PROBLEMA 3b)
F_3b := VectorField(<x, y, z>):
rot_F := Curl(F_3b);

# 4. INTEGRAL DE LÍNEA Y TRAYECTORIA POLIGONAL (PROBLEMA 4)
F_4 := VectorField(<4*x*exp(z), cos(y), 2*x^2*exp(z)>):
f_potencial := ScalarPotential(F_4);

P0 := [0, 0, 0]:
P1 := [1, 1, 0]:
P2 := [1, 1, 1]:
P3 := [0, 1, 1]:

val_P3 := eval(f_potencial, [x = P3[1], y = P3[2], z = P3[3]]):
val_P0 := eval(f_potencial, [x = P0[1], y = P0[2], z = P0[3]]):
integral_linea := simplify(val_P3 - val_P0);

c1 := spacecurve([t, t, 0], t = 0 .. 1, color = red, thickness = 3):
c2 := spacecurve([1, 1, t], t = 0 .. 1, color = green, thickness = 3):
c3 := spacecurve([1-t, 1, 1], t = 0 .. 1, color = blue, thickness = 3):

nodos := pointplot3d([P0, P1, P2, P3], symbol = solidbox, symbolsize = 15, color = black):
etiquetas := textplot3d([[P0[1],P0[2],P0[3]-0.1, "P0(0,0,0)"], 
                         [P1[1]+0.05,P1[2],P1[3], "P1(1,1,0)"], 
                         [P2[1],P2[2],P2[3]+0.1, "P2(1,1,1)"], 
                         [P3[1]-0.1,P3[2],P3[3], "P3(0,1,1)"]], 
                        font = [HELVETICA, BOLD, 12]):

display([c1, c2, c3, nodos, etiquetas], 
        axes = normal, 
        gridlines = true,
        labels = ["X", "Y", "Z"],
        title = "Trayectoria Poligonal C en el Espacio");
\end{lstlisting}

\end{document}
